\documentclass[../../../thesis.tex]{subfiles}
\begin{document}
\subsection{Článok v odbornom periodiku}
Záznam musí obsahovať mená autorov (tvorcov), názov článku, názov časopisu, dátum alebo iba rok vydania, ročník alebo zväzok (volume) a číslo vo zväzku (issue), číslo prvej strany, prípadne rozsah strán, na ktorých sa nachádza.

\subsubsection*{\normalsize Nepovinné údaje}
Ak je článok dostupný online, je dobré, ak sa v zápise objaví aj webová adresa digitálnej verzie článku. Vedecké publikácie majú pridelený jedinečný kód, tzv. DOI (Digital Object Identifier – identifikátor digitálneho objektu), ktorý tiež môže byť súčasťou bibliografického zápisu. Ďalšou nepovinnou časťou je International Standard Serial Number (ISSN), teda medzinárodné štandardné sériové číslo, ktoré sa prideľuje časopisom.

\subsubsection*{\normalsize Printový časopis}
\begin{trivlist}
\item TVORCOVIA. Názov článku. \textit{Názov periodika.} Rok vydania, \textbf{zväzok}(časť), rozsah strán. Dostupnosť.
\end{trivlist}

\subsubsection*{\normalsize Online časopis}
\begin{trivlist}
\item TVORCOVIA. Názov článku. \textit{Názov periodika} [online]. Rok vydania, \textbf{zväzok}(časť), rozsah strán (ak je k dispozícii). [cit. yyyy-mm-dd]. Dostupnosť.
\end{trivlist}

\subsubsection*{\normalsize Príklady}
Prvý príklad je z~časti
\ref{sec:citExample} odkazuje na článok s názvom Princípy formovania vzdelania v informačnej vede od autorky Jely Steinerovej, ktorý vyšiel v roku 2000 v 3. čísle 2. ročníka časopisu Pedagogická revue. Článok sa v časopise nachádza na stranách 8 až 16.

V druhom príklade skratka et. al. (skratka latinského slovného spojenia \textit{et alii}, vo význame a~iní) za menom autora znamená, že článok má viacero autorov (bežne viac než 5) a nemusíme ich všetkých uvádzať.
V slovenských textoch sa môže použiť skratka a~kol. (a~kolektív).

Tretí príklad je príklad online časopisu, ktorý možno nájsť na internete. Príznak [online] a~dátum citovania [cit. 2023-01-10] sa objaví v zozname literatúry, keď v súbore \verb|bibliography.bib| použijeme pole \verb|urldate|.

\begin{enumerate}
\item STEINEROVÁ, J. Princípy formovania vzdelania v informačnej vede. \textit{Pedagogická revue}. 2000, \textbf{2}(3), 8–16.
\item BEŇAČKA, J. et al. A better cosine approximate solution to pendulum equation. \textit{International Journal of Mathematical Education in Science and Technology}, 2009, \textbf{40}(2), 206–215.
\item HOGGAN, D. B. Challenges, Strategies, and Tools for Research Scientists. Electronic Journal of \textit{Academic and Special Librarianship} [online]. 2009, \textbf{3}(3). [cit. 2023-01-10]. Dostupné z: \texttt{http://southernlibrarianship.icaap.org/content/ v03n03/Hoggan\_d01.htm}. ISSN 1525-321X.
\end{enumerate}

\subsubsection*{\normalsize Zápis v zdrojovom .bib súbore}
\begin{verbatim}
@article{Steinerova2000Principy,
  author  = {J. Steinerová},
  journal = {Pedagogická revue},
  title   = {Princípy formovania vzdelania v informačnej vede},
  year    = {2000},
  number  = {3},
  pages   = {8--16},
  volume  = {2},
}

@article{Benacka2009Abetter,
  author  = {Ján Benačka and others},
  title   = {A better cosine approximate solution to pendulum
             equation},
  journal = {International Journal of Mathematical Education in Science
             and Technology},
  volume  = {40},
  number  = {2},
  pages   = {307--308},
  year    = {2009},
  doi     = {10.1080/00207390802419594},
}

@article{Hoggan2002Challenges,
  author  = {Danielle Bodrero Hoggan},
  journal = {Electronic Journal of Academic and Special Librarianship},
  title   = {Challenges, Strategies, and Tools for Research Scientists:
             Using Web-Based Information Resources},
  year    = {2002},
  number  = {3},
  volume  = {3},
  url     = {https://southernlibrarianship.icaap.org/content/v03n03/
             Hoggan_d01.htm},
  urldate = {2023-01-10},
}
\end{verbatim}
\end{document}